\section{The ThinkingTrees Package}\label{sec:software}

We release \textsc{ThinkingTrees}, an open-source Python package implementing the full Oracle-Preserving Summarization pipeline. The package comprises approximately 182,000 lines of Python code with comprehensive test coverage, designed for both research prototyping and production deployment. This section describes the architecture; Appendix~\ref{app:software} provides installation instructions and API documentation.

\subsection{Architecture Overview}

The package follows a plugin-based architecture with four main layers:

\begin{description}[nosep]
    \item[Core Layer] (\texttt{src/core/}): Generic building blocks including data models, LLM clients, scoring primitives, and batch processing.
    \item[Tree Layer] (\texttt{src/tree/}): Hierarchical tree construction (\texttt{TreeBuilder}) and probabilistic verification (\texttt{TreeAuditor}).
    \item[Training Layer] (\texttt{src/training/}): DSPy optimization, preference learning, and judge implementations.
    \item[Task Layer] (\texttt{src/tasks/}): Task-specific plugins (e.g., Manifesto RILE scoring) with rubrics and oracles.
\end{description}

\subsection{ThinkingTrees Platform and \texttt{treepo} Method Package}

We keep \textsc{ThinkingTrees} as the full platform name for the end-to-end OPS system, but we now expose a focused PyTorch method package, \texttt{treepo}, for simulation-heavy work. This split is practical rather than conceptual: \textsc{ThinkingTrees} owns the broader long-document pipeline, while \texttt{treepo} owns reproducible benchmark harnesses, exact sketch sanity checks, and paper-facing simulation reports.

For the commutative-sketch bridge, \texttt{treepo} exposes a public HyperLogLog surface (\texttt{HLLConfig}, \texttt{HyperLogLogSketch}) with streaming updates, exact merges, and tree-reduction helpers. The same package also exposes two benchmark families used in our running examples: (i) end-to-end cardinality recovery, where a learned TreePO sketch is compared against exact-set truth, exact HLL, and a deliberately wrong non-merge-safe baseline; and (ii) HLL merge-learning, where the learned object is the merge law itself over exact HLL register states.

\subsection{Core Components}

\paragraph{Data Models.}
The fundamental data structure is the \texttt{Node} class, which represents a node in the hierarchical summary tree:

\begin{lstlisting}
@dataclass
class Node:
    content: str           # The summary text at this node
    children: list[Node]   # Child nodes (empty for leaves)
    metadata: NodeMetadata # Audit status, scores, etc.
\end{lstlisting}

A \texttt{Tree} is the root \texttt{Node} plus methods for traversal, serialization, and audit access. The \texttt{AuditResult} dataclass tracks violation counts and confidence intervals for each condition.

\paragraph{TreeBuilder.}
The \texttt{TreeBuilder} class constructs hierarchical summaries from documents:

\begin{enumerate}[nosep]
    \item \textbf{Chunking}: Partition the document into context-fitting blocks using configurable tokenizers and overlap settings.
    \item \textbf{Leaf summarization}: Apply the summarizer $g$ to each block, creating leaf nodes.
    \item \textbf{Recursive merging}: Build the tree bottom-up by concatenating sibling summaries and applying $g$ again.
    \item \textbf{Optional cleanup}: Apply additional rounds of summarization to the root.
\end{enumerate}

The builder supports multiple \texttt{SummarizationStrategy} backends:
\begin{itemize}[nosep]
    \item \texttt{BatchedStrategy}: High-throughput async batched inference with request pooling
    \item \texttt{DSPyStrategy}: Optimizable summarization with DSPy modules
    \item \texttt{TournamentStrategy}: Tournament-based selection with learning
\end{itemize}

\paragraph{TreeAuditor.}
The \texttt{TreeAuditor} class implements the probabilistic audit from Section~\ref{sec:audit}:

\begin{lstlisting}
auditor = TreeAuditor(oracle=oracle, confidence=0.95)
result = auditor.audit(tree, n_leaves=100, n_merges=50)
print(f"Sufficiency rate: {result.p_suff:.4f} +/- {result.ci_suff:.4f}")
\end{lstlisting}

The auditor uses Hoeffding bounds to compute confidence intervals, supports streaming mode for incremental sampling, and flags nodes exceeding violation thresholds for human review.

\subsection{DSPy-First Optimization}

The package integrates deeply with DSPy \citep{khattab2024dspy} for prompt optimization. The training layer provides:

\begin{description}[nosep]
    \item[Bootstrap Optimizer] (\texttt{optimization/bootstrap.py}): Random search over prompt variations using successful traces as exemplars.
    \item[GEPA] (\texttt{optimization/gepa.py}): Gradient-Ensemble based Prompt Augmentation.
    \item[MIPRO] (\texttt{optimization/mipro.py}): Multi-task Instruction Prompt Optimization.
\end{description}

The consistency bootstrap loop (Section~\ref{sec:audit}) is implemented as:

\begin{lstlisting}
trainer = ConsistencyBootstrap(
    leaf_module=LeafSummarizer(),
    merge_module=MergeSummarizer(),
    oracle=oracle,
    optimizer="mipro"
)
trainer.train(documents, target_violation_rate=0.01)
\end{lstlisting}

\subsection{Preference Learning Integration}

For downstream preference learning, the package provides:

\begin{itemize}[nosep]
    \item \textbf{PreferencePair}: Abstract representation of (chosen, rejected) pairs with oracle scores.
    \item \textbf{PreferenceCollector}: Generate pairs from oracle comparisons on summaries.
    \item \textbf{Judge implementations}: Oracle-based, LLM-based, and GenRM-based pairwise judges.
    \item \textbf{TRL integration}: Export to Hugging Face TRL format for DPO/PPO training.
\end{itemize}

The key insight from Theorem~\ref{thm:dpo-equiv} is that when the audit certifies small violation rates, preferences collected from summaries are equivalent to preferences from full documents.

\subsection{Task Plugin Architecture}

Tasks are defined via the \texttt{AbstractTask} interface:

\begin{lstlisting}
class AbstractTask(Protocol):
    def oracle(self, text: str) -> OracleScore: ...
    def rubric(self) -> str: ...
    def metric(self, a: OracleScore, b: OracleScore) -> float: ...
\end{lstlisting}

The package includes two complete task implementations:

\begin{description}[nosep]
    \item[Manifesto RILE] (\texttt{tasks/manifesto/}): Left-right scoring of party manifestos on a $-100$ to $+100$ scale. Includes expert-calibrated rubrics, data loaders for the Manifesto Project corpus, and comparison metrics.
    \item[Congressional Record] (\texttt{tasks/congressional/}): Policy position extraction on specific issues. (In development; see Section~\ref{sec:congressional}.)
\end{description}

New tasks are registered via the plugin system:

\begin{lstlisting}
@register_task("my_custom_task")
class MyTask(ScoringTask):
    scale = BoundedScale(min_val=-100, max_val=100)
    rubric = "Extract the author's position on..."
\end{lstlisting}

\subsection{LLM Backend Support}

The package supports multiple LLM backends through a unified \texttt{LLMClient} interface:

\begin{itemize}[nosep]
    \item \textbf{OpenAI API}: GPT-4, GPT-4o, o1, etc.
    \item \textbf{vLLM}: Self-hosted inference with batching
    \item \textbf{SGLang}: High-performance serving
\end{itemize}

The \texttt{AsyncBatchLLMClient} provides request pooling and rate limiting for high-throughput workloads.

\subsection{Configuration}

Configuration is YAML-based with sensible defaults:

\begin{lstlisting}
model:
  name: gpt-4o
  temperature: 0.0
  max_tokens: 2048
chunking:
  max_chunk_size: 4000
  overlap: 200
audit:
  n_leaves: 100
  n_merges: 50
  confidence: 0.95
\end{lstlisting}

The full configuration schema and options are documented in Appendix~\ref{app:software}.
